\chapter{Análisis de riesgos}
\label{cap:riesgos}
En este capítulo se realizará un análisis de los posibles riesgos que pueden afectar al proyecto,
así como de las medidas que se pueden tomar para mitigarlos. Se clasificarán según la probabilidad
de que ocurran y el impacto que tendrían en el sistema, ambos valores enteros entre uno y cinco.
También se dará la magnitud de exposición al riesgo, en este caso definida como el producto de
la probabilidad y el impacto \parencite[p.~56]{MariaMedicion}.

% Probabilidad: Muy improbable (1), improbable (2), Posible (3), Probable (4), Muy probable (5)
% Impacto: Inegligible (1), Bajo (2), Moderado (3), Alto (4), Catastrófico (5)

\riesgo{
    id = {R1},
    nombre = {Baja adopción},
    descripcion = {%
        Aunque el sistema se haya diseñado para eliminar las barreras de entrada a los \glspl{exergame},
        el público puede ser reacio a adoptar esta nueva forma de hacer ejercicio. Es posible que la
        novedad no sea suficiente para atraer a los usuarios.
    },
    probabilidad = {Posible (3)},
    impacto = {Alto (4)},
    exposicion = {Riesgo medio (12)},
}

% Insatisfacción por imprecisión y latencia 4 3
\riesgo{
    id = {R2},
    nombre = {Insatisfacción por imprecisión y latencia},
    descripcion = {%
        La calidad de la experiencia dada por el producto creado depende en gran medida de la detección
        de posturas en tiempo real a través de la cámara. A pesar de la capacidad de adaptar la calidad
        del seguimiento, los dispositivos de gamas más bajas verán reducidos el rendimiento y/o la
        calidad de la detección, afectando a la jugabilidad.
    },
    probabilidad = {Probable (4)},
    impacto = {Moderado (3)},
    exposicion = {Riesgo medio (12)},
}

% Falta de profundidad y rejugabilidad 4 4
\riesgo{
    id = {R3},
    nombre = {Falta de profundidad y rejugabilidad},
    descripcion = {%
        El producto se lanza únicamente con dos minijuegos posibles. Estos dos minijuegos son
        experiencias completas, que permiten escoger diferentes niveles de dificultad, aumentando
        la rejugabilidad y el tiempo de vida de la aplicación. Sin embargo, es probable que, tras
        poco tiempo, los jugadores pierdan el interés por las pocas opciones disponibles, afectando
        al uso de la plataforma.
    },
    probabilidad = {Probable (4)},
    impacto = {Alto (4)},
    exposicion = {Riesgo alto (16)},
}

% Dependencia del entorno 5 3
\riesgo{
    id = {R4},
    nombre = {Dependencia del entorno},
    descripcion = {%
        La eficacia del modelo de detección de posturas depende en gran medida de factores externos
        como la calidad de la iluminación o la complejidad del fondo de la escena. Estas condiciones,
        bastante comunes en entornos domésticos, pueden afectar a la fiabilidad de la detección,
        haciendo que la aplicación sea inutilizable para los usuarios.
    },
    probabilidad = {Muy probable (5)},
    impacto = {Moderado (3)},
    exposicion = {Riesgo alto (15)},
}

% Problemas de compatibilidad 2 4
\riesgo{
    id = {R5},
    nombre = {Problemas de compatibilidad},
    descripcion = {%
        El proyecto se ejecuta en navegadores utilizando \texttt{WebGL}, \texttt{WebAssembly} y
        \texttt{MediaDevices}. Las implementaciones de estas tecnologías pueden variar
        significativamente entre navegadores web y sistemas operativos, lo que podría llevar a
        errores o a caídas de rendimiento imprevistas.
    },
    probabilidad = {Improbable (2)},
    impacto = {Alto (4)},
    exposicion = {Riesgo medio (8)},
}

% Obsolescencia tecnológica 3 2
\riesgo{
    id = {R6},
    nombre = {Obsolescencia tecnológica},
    descripcion = {%
        Existe el riesgo de que las tecnologías utilizadas se vuelvan obsoletas, haciendo la labor
        de mantenimiento más compleja.
    },
    probabilidad = {Posible (3)},
    impacto = {Bajo (2)},
    exposicion = {Riesgo medio (6)},
}

